\sscp{Major Austronesian incursion: west-to-east along southern islands}{15-06-01}
There are retentions of many features in the historical phonology,
morphosyntax, and lexicon in the \tsc{Bima-Lembata} subgroup
(\crf{ch:BL}) that are fairly conservative,
and are often associated more with languages west of the Blust
line than with those to the east.
Furthermore, many of the features resulting
from contact identified in \srf{15-Features} and \srf{15-04},
are not found or are found in very limited ways in western,
central and southern parts of the \tsc{Bima-Lembata} subgroup.
It is mostly in the eastern parts of the subgroup that many of
these contact-related features start to be found more systematically.
In particular, the \tsc{Flores-Lembata} node of \tsc{Bima-Lembata}
shows the most signs of contact-induced change \citep{Fricke2019}.
Innovations define subgroups; retentions do not define subgroups.
But retentions also tell an important part of the fuller story.

Together these features imply a scenario something like moving
from the west, perhaps across Java, Bali,
Lombok, Sumbawa, Flores southwardswards to Sumba to Sabu,
and Flores eastwards to Adonara to Lembata.
Distances between islands are relatively short.
For example, the straits between Flores
and Adonara is much narrower than many large rivers of the world.
It is not yet clear from the linguistic evidence if the Austronesian-speaking
peoples continued on and reached Timor from this incursion,
or stopped at Pantar and Alor.
Evidence for this scenario includes:

\begin{itemize}
	\item	Retention of PMP *ə = \it{ə}: This is strong
				and widespread, and is evident in many six-vowel systems across
				the whole subgroup, but not necessarily in all languages.
				Other Wallacean subgroups have five-vowel systems.
				See \srf{03-02-01} and \srf{03-03-01}.
	\item	Retention of vowel-glide sequences: *-ay > \it{aʤ,
				e, a, i}; *-aw > \it{av,
				o, a,} \it{u}; *-uy > \it{u, i}.
				These on-the-ground forms require reconstruction of the full sequences.
	\item	Retention of *b > \it{b,
				w}; *d = \it{d} \citep[56, 85f]{Blust2008}.
	\item	Retention of *p = \it{p}.
	\item	Retention of *ŋ = \it{ŋ}.
				This is strong and widespread.
	\item	Partial retention of *ñ = \it{ɲ} \citep[57, 68]{Blust2008}.
	\item	Retention of many medial *NC clusters \citep[54]{Blust2008}. See \srf{13-02-01}.
	\item	Retention of medial *q \citep[76f, 85]{Blust2008}.
	\item	Partial retention of intervocalic *y \citep[88]{Blust2008}.
	\item	Retention of productive PMP verb morphology *ma- stative,
				*pa- causative, and P(CE)MP *pa(Ra)- reciprocal (reflexes of
				PWMP *paR- also have ‘reciprocal’ as one of its many functions.
				See \srf{13-03-01}.) 
\end{itemize}

A number of words more commonly associated with the CEMP region
are not so commonly found in much of the \tsc{Bima-Lembata} subgroup (see \srf{15-04}).
It is fair to ask if words associated with languages west of
the Blust line are retentions from an incursion from the west,
or relexification due to later contact with languages west of the Blust line.
The collective patterns in the historical phonology,
regular sound correspondences in the lexicon (loan words often
do not follow regular sound correspondences),
and other patterns of contact favour retention over later borrowing
from western languages as the more likely explanation.

As mentioned above, many structural features in the wider region
attributed to contact with SOV Papuan languages are much less common in \tsc{Bima-Lembata}.
They are more common in the eastern parts
and less common in the western
and southern parts of the subgroup.
For example, post-posed possessors are the dominant pattern in
western and southern regions of \tsc{Bima-Lembata},
with optional post-posed or pre-posed constructions found in
the eastern edge of the subgroup.
Clause-final negation and alienable-inalienable distinctions
also begin to be found in the eastern parts of the subgroup.
Other features associated with Papuan languages such as {\#}muku
‘banana’ reach as far west as Manggarai in western Flores
(\srf{02-03-01-07}).
There is a strong bundle of evidence from the historical phonology,
the lexicon, the morphology,
and syntactic structures that a group of people speaking Austronesian
languages not very divergent from PMP entered the \tsc{Bima-Lembata}
region probably in the west.
They travelled west to east,
and some went south to Sumba and Sabu.
There are more features ascribed to contact-induced change in
the eastern parts of the subgroup,
and fewer in the western regions,
suggesting more contact, more intense contact
and more bilingualism the further east they travelled.

But drawing a clear link in the historical phonology between
\tsc{Bima-Lembata} and \tsc{Timor-Babar} below the level of PMP is difficult.
Or to put it in unscientific terms,
“It looks like they got off a different boat.” These discontinuities
are reflected in several distributional maps throughout this study.